\chapter{Gonorrhea}
\index{Gonorrhea}
\label{sec:Gonorrhea}

\section{Overview}
Gonorrhea is a sexually transmitted infection caused by the gram-negative diplococcus \textit{Neisseria gonorrhoeae}. It primarily infects the mucous membranes of the urogenital tract, rectum, pharynx, and conjunctivae. Untreated gonorrhea can lead to serious reproductive complications and facilitate HIV transmission.
\section{Classification}

\begin{description}[font=\normalfont\bfseries,leftmargin=3em,labelwidth=2.5em]
  \item[Cause] Infectious (bacterial)
  \item[Organ System] Infectious/Reproductive
  \item[Pathophysiology] \textit{N. gonorrhoeae} attaches to columnar epithelial cells via pili and Opa proteins, triggers local inflammation with neutrophil infiltration, and can evade immune clearance through antigenic variation
  \item[Duration] Acute (weeks if treated); chronic if untreated
  \item[Onset] Sudden (2--14 days after exposure)
  \item[Mode of Transmission] Sexual contact (vaginal, anal, oral); perinatal
  \item[Clinical Course] Acute urethritis or cervicitis with purulent discharge; may progress to ascending infection or become asymptomatic
  \item[Severity] Mild to moderate; severe with disseminated disease
  \item[Extent of Disease] Localized (urogenital, rectal, pharyngeal); systemic in disseminated gonococcal infection
  \item[Age Group] Sexually active adolescents and young adults (15--29 years)
  \item[Epidemiology] Second most commonly reported bacterial STI worldwide; highest incidence in 15--24-year-old women and 20--29-year-old men
  \item[ICD-11 Code] AA00.0
\end{description}

\section{Causes}
Gonorrhea is caused by infection with \textit{Neisseria gonorrhoeae}, an aerobic gram-negative bacterium that colonizes mucosal surfaces. The bacterium penetrates columnar and transitional epithelium, triggering an intense neutrophilic inflammatory response that produces the characteristic purulent discharge.

\section{Risk Factors}

\begin{itemize}[noitemsep]
  \item Unprotected vaginal, anal, or oral sex
  \item Multiple sexual partners or a partner with multiple partners
  \item Age 15--29 years
  \item Previous or concurrent STI (especially chlamydia)
  \item Inconsistent condom use
  \item Sex work or transactional sex
  \item Substance use associated with high-risk sexual behavior
  \item Men who have sex with men (MSM)
\end{itemize}

\section{Signs and Symptoms}

\subsection{Early symptoms}
\begin{itemize}[noitemsep]
  \item Early manifestations of gonorrhea may be subtle and nonspecific
\end{itemize}

\subsection{Common symptoms}
\begin{itemize}[noitemsep]
  \item \subsection{Common symptoms}
  \item \textbf{Common symptoms:}
  \item Purulent urethral discharge (men)
  \item Dysuria and urinary urgency
  \item Cervicitis with mucopurulent discharge (women)
  \item Intermenstrual bleeding or menorrhagia
  \item Anal discharge, itching, or pain (rectal infection)
  \item Sore throat (pharyngeal infection, often asymptomatic)
  \item Conjunctivitis with purulent discharge (ocular infection)
  \item Testicular pain or swelling (epididymitis)
  \item Lower abdominal or pelvic pain
  \item Many infections (especially in women) are asymptomatic
\end{itemize}

\subsection{Less common symptoms}
\begin{itemize}[noitemsep]
  \item Atypical or less common presentations of gonorrhea may occur
\end{itemize}

\subsection{Emergency warning signs}
\begin{itemize}[noitemsep]
  \item Severe or worsening symptoms of gonorrhea require immediate medical evaluation
\end{itemize}
\section{Progression and Stages}
Incubation is typically 2--8 days. Acute infection presents with urethritis or cervicitis; without treatment, the infection can ascend to cause epididymitis in men and pelvic inflammatory disease (PID) in women. Hematogenous dissemination occurs in 1--3\% of cases, causing dermatitis, tenosynovitis, and septic arthritis.

\section{Complications}

\begin{itemize}[noitemsep]
  \item Pelvic inflammatory disease (PID) with tubal scarring
  \item Infertility (tubal factor in women; epididymal obstruction in men)
  \item Ectopic pregnancy
  \item Disseminated gonococcal infection (DGI) with dermatitis and arthritis
  \item Neonatal ophthalmia (ophthalmia neonatorum) from perinatal exposure
  \item Increased HIV acquisition and transmission risk
  \item Fitz-Hugh--Curtis syndrome (perihepatitis)
  \item Reduced quality of life
  \item Emotional and psychological effects
\end{itemize}

\section{Diagnosis}

\begin{itemize}[noitemsep]
  \item Nucleic acid amplification test (NAAT) of urine, urethral, cervical, rectal, or pharyngeal swab
  \item Gram stain of urethral discharge showing gram-negative intracellular diplococci (men)
  \item Culture with antimicrobial susceptibility testing (if treatment failure suspected)
  \item Test for concurrent chlamydia, HIV, and syphilis
  \item Pelvic examination in women with suspected PID
  \item Lab tests
  \item Imaging tests (pelvic ultrasound for suspected PID)
\end{itemize}

\section{Prevention}

\begin{itemize}[noitemsep]
  \item Consistent and correct condom use
  \item Routine STI screening for sexually active young women and MSM
  \item Partner notification and treatment (expedited partner therapy)
  \item Avoid sexual contact until treatment is complete
  \item Regular check-ups
\end{itemize}

\section{Treatment Options}

\begin{itemize}[noitemsep]
  \item Ceftriaxone 500 mg intramuscularly as a single dose (first-line)
  \item Treatment for concurrent chlamydia (doxycycline 100 mg twice daily for 7 days)
  \item Test of cure for pharyngeal infections
  \item Sexual partners should be treated empirically
  \item Lifestyle changes
\end{itemize}

\section{Living with It}

\begin{itemize}[noitemsep]
  \item Follow treatment plan
  \item Regular appointments
  \item Adjust daily habits
  \item Support groups
  \item Keep learning
\end{itemize}

\section{Outlook}
With appropriate antibiotic treatment, uncomplicated gonorrhea is completely curable and symptoms resolve within days. However, reinfection is common without partner treatment and behavior change. Antimicrobial resistance (especially extended-spectrum cephalosporin resistance) is a growing global concern that may complicate future treatment.
